\documentclass[12pt,a4paper]{article}

\usepackage[spanish]{babel}
\usepackage[utf8]{inputenc}
\usepackage[T1]{fontenc}
\usepackage{geometry}
\usepackage{setspace}
\usepackage{titlesec}
\usepackage{graphicx}
\usepackage{tikz}
\usetikzlibrary{shapes.geometric, arrows.meta, babel} % Agregamos shapes y babel
\usepackage{enumitem}
\usepackage{amsmath}

\geometry{margin=2.5cm}
\setstretch{1.5}

\title{\textbf{Modelamiento de Procesos con BPMN}\\
\large Enfoque en Ingeniería de Sistemas}
\author{}
\date{}

\begin{document}

\maketitle

\tableofcontents
\newpage

\section{Concepto de BPMN}

\subsection{Definición}

BPMN (Business Process Model and Notation) es un estándar internacional de modelamiento gráfico que permite representar procesos de negocio de forma clara, formal y comprensible tanto para usuarios técnicos como no técnicos.

Desde una perspectiva académica, BPMN constituye una notación visual que describe la lógica de los procesos organizacionales mediante símbolos estandarizados que representan actividades, eventos, decisiones y flujos.

\subsection{Aplicación en las organizaciones}

El uso de BPMN en las organizaciones permite:

\begin{itemize}
    \item Documentar procesos empresariales
    \item Analizar y mejorar flujos de trabajo
    \item Facilitar la comunicación entre áreas de negocio y tecnología
    \item Automatizar procesos mediante sistemas BPM
\end{itemize}

En el contexto de la ingeniería de sistemas, BPMN permite traducir procesos organizacionales en modelos estructurados susceptibles de implementación en software.

\subsection{Importancia en el modelamiento de procesos}

La relevancia de BPMN radica en:

\begin{itemize}
    \item Proporcionar un lenguaje común organizacional
    \item Identificar ineficiencias y redundancias
    \item Facilitar la estandarización de procesos
    \item Servir como base para la automatización
\end{itemize}

Esto lo convierte en una herramienta clave en la transformación digital y la gestión por procesos.

\section{Reglas generales para el modelamiento en BPMN}

\subsection{Uso de pools}

Un \textit{pool} representa un participante principal dentro del proceso, como una organización o sistema.

Se utiliza cuando existen múltiples actores independientes o cuando se requiere modelar interacción entre entidades.

\subsection{Uso de lanes}

Los \textit{lanes} son subdivisiones dentro de un pool que representan roles, áreas o departamentos. Permiten organizar responsabilidades y visualizar quién ejecuta cada actividad.

\subsection{Niveles de detalle}

El modelamiento BPMN puede realizarse en distintos niveles:

\begin{itemize}
    \item Alto nivel: visión general del proceso
    \item Nivel detallado: incluye reglas, decisiones y tareas específicas
\end{itemize}

\subsection{Orientación del diagrama}

Los diagramas deben seguir una dirección lógica:

\begin{itemize}
    \item De izquierda a derecha
    \item De arriba hacia abajo
\end{itemize}

\subsection{Buenas prácticas}

\begin{itemize}
    \item Claridad y simplicidad
    \item Evitar cruces innecesarios
    \item Uso de nombres descriptivos
    \item Coherencia en el nivel de detalle
    \item Evitar sobrecarga visual
\end{itemize}

\section{Elementos BPMN}

\subsection{Conectores}

\subsubsection{Flujo de secuencia}

Representa el orden de ejecución de actividades dentro de un mismo proceso.

\subsubsection{Flujo de mensaje}

Indica comunicación entre diferentes participantes (pools).

\subsubsection{Asociación}

Relaciona elementos adicionales como datos o anotaciones sin afectar el flujo.

\subsection{Actividades}

Las actividades representan el trabajo dentro del proceso:

\begin{itemize}
    \item Actividad de usuario
    \item Actividad manual
    \item Script task
    \item Service task
    \item Send task
    \item Receive task
\end{itemize}

\subsection{Subprocesos}

\subsubsection{Subproceso embebido}

Forma parte del proceso principal y simplifica la visualización.

\subsubsection{Subproceso reutilizable}

Es independiente y puede ser utilizado en múltiples procesos.

\subsection{Eventos}

\begin{itemize}
    \item Evento de inicio
    \item Evento intermedio
    \item Evento de fin
\end{itemize}

\subsection{Compuertas (Gateways)}

\begin{itemize}
    \item Exclusiva (XOR)
    \item Paralela (AND)
    \item Inclusiva (OR)
    \item Basada en eventos
\end{itemize}

\section{Caso práctico: Proceso de compra en línea}

\subsection{Descripción del proceso}

Se modela el proceso desde que un cliente realiza una compra hasta la entrega del producto.

\subsection{Estructura organizacional}

\begin{itemize}
    \item Pool Cliente
    \item Pool Empresa
    \begin{itemize}
        \item Ventas
        \item Logística
        \item Finanzas
    \end{itemize}
\end{itemize}

\subsection{Diagrama del proceso}

\begin{center}
\begin{tikzpicture}[
    node distance=1.8cm, 
    every node/.style={draw, rectangle, rounded corners, align=center, minimum width=3cm}
]

% Nodos alineados verticalmente
\node (start) {Inicio};
\node (pedido) [below of=start] {Realizar pedido};
\node (validar) [below of=pedido] {Validar datos};
\node (decision1) [below of=validar, diamond, aspect=2, minimum width=2cm] {¿Datos válidos?};
\node (pago) [below of=decision1, yshift=-0.5cm] {Procesar pago};
\node (decision2) [below of=pago, diamond, aspect=2, minimum width=2cm] {¿Pago aprobado?};
\node (paralelo) [below of=decision2, yshift=-0.5cm] {Procesos paralelos};
\node (envio) [below of=paralelo] {Despachar producto};
\node (fin) [below of=envio] {Fin};

% Flechas
\draw[->] (start) -- (pedido);
\draw[->] (pedido) -- (validar);
\draw[->] (validar) -- (decision1);
\draw[->] (decision1) -- node[right]{Sí} (pago);
\draw[->] (pago) -- (decision2);
\draw[->] (decision2) -- node[right]{Sí} (paralelo);
\draw[->] (paralelo) -- (envio);
\draw[->] (envio) -- (fin);

% Opcional: Caminos de salida para "No" (ejemplo)
% \draw[->] (decision1) -- ++(-3,0) node[above]{No} |- (pedido);

\end{tikzpicture}
\end{center}

\subsection{Flujo del proceso}

El proceso incluye:

\begin{enumerate}
    \item Inicio con solicitud del cliente
    \item Registro de datos
    \item Validación automática
    \item Evaluación mediante compuertas XOR
    \item Procesamiento de pago
    \item Ejecución paralela de tareas
    \item Envío del producto
    \item Finalización del proceso
\end{enumerate}

\subsection{Interpretación}

El modelo evidencia:

\begin{itemize}
    \item Interacción entre actores mediante flujo de mensajes
    \item Toma de decisiones con compuertas XOR
    \item Ejecución simultánea con compuertas AND
    \item Integración de diferentes tipos de actividades
\end{itemize}

\section{Conclusión}

BPMN constituye una herramienta fundamental en la ingeniería de procesos y sistemas, ya que permite representar de manera estructurada y estandarizada los procesos organizacionales.

Su correcta aplicación facilita la comprensión, análisis y mejora de los flujos de trabajo, además de servir como base para la automatización de procesos.

El dominio de sus elementos —eventos, actividades, compuertas, conectores y subprocesos— permite construir modelos precisos que reflejan la realidad organizacional, contribuyendo significativamente al diseño de sistemas de información y a la optimización empresarial.

\end{document}
```

